\begin{lemma}[Stable change, Cases 1 and 2, with a transported starred tail]
Fix a finite loopless multigraph \(G\) with its individual edge type, a
palette of \(k\) colors, a distinguished edge \(e_0\), an input coloring,
an ordered tree-sequence \(T\), a generating coloring \(\varphi_n\), a
positive-rung ETT \(\mathcal E\) on \(T\), a hierarchy base, and a
\noderef{HierarchyCounterexample} \(C\).  Assume the properness,
exactly-one-uncoloredness, complete lower-rung six-pack, fewer-divider
induction, rung-local lexicographic minimality, and \(2n+11\) estimate
recorded in a
\noderef{HierarchyMinimumCounterexampleContext}.

Write \(H=C.\mathit{hierarchy}\), let \(q\) be its last stored index, and
put
\[
 A=\noderef{HierarchyStarredStage}(H,q),\qquad
 B=\noderef{HierarchyFullTailPrefix}
       (C,\noderef{HierarchyFullTailLength}(C)-1),
\]
\[
 D_{n,q}=\noderef{HierarchyStoredColorsAt}(H,q),\qquad
 \Gamma^q=\bigcup_{\eta\in D_{n,q}}C.\mathit{good}.\Gamma^q_\eta .
\]
Let \(\alpha\ne\beta\) be missing on \(B\).  Let \(z\) certify a possibly
singleton \noderef{TwoColorPathComponent} \(Q\) under \(\varphi_n\), put
\[
 \sigma_n=\noderef{TwoColorComponentSwap}
      (G,\varphi_n,\alpha,\beta,z),
\]
and assume that \(Q\) is disjoint from the full
\noderef{TwoColorComponent} through the
first displayed vertex of \(T\) missing \(\alpha\), where that position is
given by \noderef{HierarchyFirstMissingPosition}.  Assume the explicit
Cases 1/2 alternative: either \(q>0\) and \(\alpha\) is missing on
\(T_{n,q}\), or \(q=0\) and either \(\alpha\) is missing on \(T_n\), or
\(\alpha\) is missing on \(A\) and both \(\alpha,\beta\) lie outside
\(D_n=\noderef{TashkinovStoredColors}(\mathcal E,n)\).

Assume in addition the following starred-stage transport data.
\begin{enumerate}
\item Both \(\alpha\) and \(\beta\) belong to
\(\noderef{TreeMissingColors}(\sigma_n,A)\).
\item Both \(\alpha\) and \(\beta\) lie outside \(\Gamma^q\).
\item For every \(v\in V(B)\),
\(\noderef{MissingColors}(G,\sigma_n,v)=
 \noderef{MissingColors}(G,\varphi_n,v)\).
\item For every \(\eta\in D_{n,q}\),
\[
 \noderef{HierarchyFirstMissingPosition}(\sigma_n,T,\eta)=
 \noderef{HierarchyFirstMissingPosition}(\varphi_n,T,\eta).
\]
\end{enumerate}

Then \(\sigma_n\) is \noderef{HierarchyWeaklyStable} at \(A\) and is a
\noderef{ModColoring} of \(\mathcal E\).  Moreover there exist an actual
\noderef{CorrespondingExtendedTashkinovTree} certificate on the same full
ordered tree \(T\), \noderef{MaximumProperty} for its replacement ETT, a
hierarchy base for that ETT, a full hierarchy \(K\), and a
\noderef{GoodHierarchy} \(g_K\), such that
\[
                 \noderef{SameGoodHierarchy}(C.\mathit{good},g_K).
\]
No equality of the affected displayed edge values after \(A\) is assumed
or concluded.
\end{lemma}
\begin{proof}
Let \(T_n\) denote the terminal ladder tree, and retain the notation
\(H,q,A,B,D_{n,q},\Gamma^q,Q\) from the statement.  Apply
\noderef{HierarchyCounterexamplePenultimateElementary} to the explicit
data and, in order, the six projections
\[
 \mathit{ctx.proper},\quad \mathit{ctx.oneEdgeUncolored},\quad
 \mathit{ctx.lowerRung},\quad \mathit{ctx.fewerDividers},\quad
 \mathit{ctx.lexicographicMinimality},\quad
 \mathit{ctx.largeMissingSet}
\]
of the stated \noderef{HierarchyMinimumCounterexampleContext}.  Its
conclusion is
\[
 V(B)\text{ is elementary under }\varphi_n.                    \tag{1}
\]
Thus fact (6.6), which is not a field of the context, has been derived.
The same six projections in
\noderef{HierarchyCounterexampleTailAtLeastTwo} give
\[
 \ell=\noderef{HierarchyFullTailLength}(C)
       =|T|-|A|\geq2.                                          \tag{2}
\]
Expanding \noderef{HierarchyFullTailPrefix} as the appropriate
\noderef{TreeSequencePrefix}, (2) and the hierarchy inclusions give exact
\noderef{TreeSequenceInitialSegment}s
\[
                         A\preccurlyeq B\preccurlyeq T,
 \qquad |B|=|T|-1.                                             \tag{3}
\]

Now apply \noderef{StableChangeCasesOneTwo} to the same ambient data.
Pass the six context projections above, (1), the two colors missing on
\(B\), the path certificate for \(Q\), the defining swap equality, the
Cases 1/2 alternative, and the component-disjointness hypothesis.  Retain
its conclusions
\[
\begin{array}{ll}
\text{(S1)}&\sigma_n\text{ is weakly stable at }A,\\
\text{(S2)}&\sigma_n\text{ is a mod-coloring of }\mathcal E,\\
\text{(S3)}&\text{if }q>0,\text{ there are an input coloring, an actual
 ETT }\mathcal E_A\text{ on }A,\\
&\text{a base }\mathcal B_A,\text{ a hierarchy }H_A,\text{ and a good
 hierarchy }g_A\text{ under }\sigma_n.
\end{array}                                                     \tag{4}
\]
Retain every further conjunct of (S3): \(\mathcal E_A\) has maximum
property; its terminal extension kind equals that of \(\mathcal E\);
\(H_A\) has one fewer divider than \(H\); the core and closure of
\(\mathcal B_A\) equal those of the original base; every stage of
\(H_A\) is the corresponding retained stage of \(H\); and, at every
retained stored index, the two \noderef{HierarchyStoredColorsAt} sets
and every Gamma pair indexed by their common members agree.  Finally
retain the separately stated last clause: \(T_{n,q}\) is minus-closed
under \(\sigma_n\) for the union of the preceding Gamma pairs indexed
by \(D_{n,q}\).  We use the names
\[
 \mathcal E_A,\quad \mathcal B_A,\quad H_A,\quad g_A,
 \quad M_A,\quad \kappa_A=\kappa,
 \quad d_A+1=q                                             \tag{4a}
\]
for these actual objects, its maximum witness, its terminal-kind
equality, and its divider-count equality, respectively.  Thus (4a) is
not a summary of unnamed existence data: all the core, closure, stage,
stored-set, Gamma-pair, and last minus-closed equalities just listed are
fixed witnesses from the conclusion of
\noderef{StableChangeCasesOneTwo}.
The first two are already the first two conclusions of this lemma.  We
construct all objects in the remaining existential conclusion.

\textbf{The exact old route.}
Let \(n\) be the last positive rung.  Inspect
\(\mathcal E.\mathit{series}.\mathit{certified}(n)\) in the
\noderef{TashkinovSeries}.  Positivity of the ETT excludes the initial
constructor of \noderef{TashkinovSeriesStep}.  In the PE constructor, its
\(\varphi_n\)-\noderef{TashkinovClosure} contains a
\(\varphi_n\)-\noderef{TashkinovAugmentationSequence} from \(T_n\) to the
certified final series tree.  In the SE and RE constructors it supplies,
respectively from the exact constructor payload, an individual edge
\(f_n\), the literal singleton identity \(F_n=\{f_n\}\), an exact
\noderef{TreeSequenceAdjoinEdge} seed \(S=T_n+f_n\), and a
\(\varphi_n\)-augmentation sequence from \(S\) to that final tree.

The \texttt{belowFinal} field of \noderef{ExtendedTashkinovTree} makes
\(T\) an exact prefix of the certified final tree.  In PE its
\texttt{abovePenultimate} field gives \(T_n\preccurlyeq T\).  In SE/RE
the seed and \(T\) are exact prefixes of the same final tree, while
\[
 |S|=|T_n|+1\leq |T|
\]
by the adjoin-edge length and the \texttt{strictAbovePenultimate} field.
Comparing the entrywise prefix equalities therefore gives
\(S\preccurlyeq T\).  Applying
\noderef{TashkinovAugmentationSequenceRestrict} gives the exact old routes
\[
\begin{cases}
 T_n\longrightarrow T,&\text{in PE},\\
 S=T_n+f_n\longrightarrow T,&\text{in SE/RE}.
\end{cases}                                                     \tag{5}
\]

We locate \(A\) in these finite routes.  If \(q>0\), unfold
\noderef{HierarchyStarredStage}; the nested and strict fields of
\noderef{TashkinovHierarchy} identify \(A\) with the last nonterminal
hierarchy stage and put the appropriate source in (5) at or before \(A\).
For SE/RE, the source \(S\) has length \(|T_n|+1\), while the first strict
hierarchy inclusion puts \(|A|\geq |T_n|+1\); since both are exact prefixes
of \(T\), this also gives \(S\preccurlyeq A\).  If \(q=0\) and the kind is
PE, then \(A\) is the base core; the left-prefix field of the
\noderef{OrderedTreeUnion} in \noderef{HierarchyBase} gives
\(T_n\preccurlyeq A\), and (3) gives \(A\preccurlyeq T\).

Every nonreflexive step of (5) is a
\noderef{TashkinovAugmentation}, hence raises length by exactly one and
preserves all earlier exact entries.  Starting from its source, take the
first member whose length is at least \(|A|\).  In the cases just listed
such a member exists; the one-step length equality gives it length
\(|A|\), and two equal-length exact prefixes of \(T\) are equal entry by
entry.  Thus it is exactly \(A\).  Factoring the reflexive-transitive chain
there gives an old prefix route to \(A\) and the old suffix
\[
                         A\longrightarrow T                    \tag{6}
\]
under \(\varphi_n\).  The sole omitted case is \(q=0\) with kind SE or RE:
then the nonparallel field of \noderef{HierarchyBase} says \(A=T_n\), and
the exact non-TAA seed \(S=T_n+f_n\) in (5) must remain intact.

\textbf{Transport to the starred stage.}
First suppose \(q>0\).  Unfolding (S1) through
\noderef{HierarchyWeaklyStable} supplies
\(\noderef{TreeColoringInvariant}(A,\varphi_n,\sigma_n)\).  Consider one
old augmentation \(X\to Y\) in the prefix part of (6).  Both \(X\) and
\(Y\) are exact prefixes of \(A\).  The old missing-color witness lies at
a displayed vertex of \(X\), hence in \(V(A)\), so the vertex clause of
the invariant preserves it.  The exact new edge has position
\(|X|\geq1\): in PE the source \(T_n\) is nonempty, and in SE/RE the source
is the still larger seed.  The edge clause of the invariant therefore
preserves its color.  Length and exact-prefix clauses are independent of
the coloring.  Induction through
\noderef{TashkinovAugmentationSequence} transports the whole prefix route
to \(\sigma_n\).

Suppose next that \(q=0\) and the terminal kind is PE.  Here (S1) is the
\noderef{CombinedStableColoring} relation on \(T_n\) and the PE closure
\(R_n\).  Applying \noderef{CombinedStableClosurePreserved} shows that
every displayed edge identity of \(T_n\) and \(R_n\) has the same value
under \(\varphi_n\) and \(\sigma_n\).  The
\noderef{OrderedTreeUnion} for the core \(A=T_n\vee R_n\) says both that
every vertex of \(A\) lies in one of these two trees and that every entry
after \(T_n\) is an exact entry of \(R_n\).  For each old augmentation in
the prefix route \(T_n\to A\), the corresponding vertex clause of
\noderef{StableColoring} preserves its missing witness, and the cited
closure-preservation result fixes its exact new-edge value.  Induction
again gives a \(\sigma_n\)-augmentation route \(T_n\to A\).

In either of these cases, the two supplied missing-color hypotheses on
\(A\), the defining equality for \(\sigma_n\), and (6) are exactly the
hypotheses of
\noderef{TashkinovAugmentationSequenceComponentSwap}.  It transports the
whole suffix, with the same ordered sequences and individual edge
identities, to a \(\sigma_n\)-augmentation sequence
\[
                         A\longrightarrow T.                  \tag{7}
\]
Concatenating the transported prefix with (7) gives the first line below
in PE and the second in positive-divider SE/RE.

Finally suppose \(q=0\) and the kind is SE or RE.  We have \(A=T_n\), and
the \noderef{TreeSequenceAdjoinEdge} certificate makes \(A\) an exact
prefix of \(S=T_n+f_n\).  Expanding \noderef{TreeMissingColors}, the two
missing witnesses on \(A\) therefore remain witnesses on \(S\).  Apply
\noderef{TashkinovAugmentationSequenceComponentSwap} directly to the old
route \(S\to T\) in (5).  This retains the same exact \(f_n\), the literal
singleton equation \(F_n=\{f_n\}\), and the same adjoin-edge seed.  We have
proved, in all cases,
\[
\begin{cases}
 T_n\longrightarrow T,&\text{under }\sigma_n\text{ in PE},\\
 T_n+f_n\longrightarrow T,&\text{under }\sigma_n\text{ in SE/RE}.
\end{cases}                                                     \tag{8}
\]
No assertion that an affected edge after \(A\) keeps its old value has
been used.

\textbf{The replacement ETT and hierarchy.}
The old ETT gives the exact strict prefix \(T_n\prec T\).  The two routes
in (8), including the singleton and seed data in the non-PE line, are
precisely the two eligibility alternatives in the universal clause of
(S2), after unfolding \noderef{ModColoring}.  Consequently (S2) produces
an actual \noderef{CorrespondingExtendedTashkinovTree} certificate
\(\mathcal C_T^{\rm mod}\) on \(T\).  Apply
\noderef{MaximumPropertyTransport} to (S2),
\(\mathcal C_T^{\rm mod}\), and \(C.\mathit{maximum}\).  This gives a
maximum witness \(M_T^{\rm mod}\) for its replacement ETT.  We retain
these two objects for the SE/RE branch.

If the terminal kind is PE, (S1) supplies the required combined stability
and the first line of (8) supplies the full route required by
\noderef{HierarchySplitter}.  The properness and exactly-one-uncolored
inputs for \(\sigma_n\) are the first two fields of (S2), while the
corresponding inputs for \(\varphi_n\) and the lower-rung package are
projections of the stated context.  Apply that lemma with these facts and
\(C.\mathit{maximum}\).  Denote the
objects returned by this single application by
\[
 \mathcal C_T,\qquad M_T,\qquad \mathcal B_T,\qquad K.          \tag{9}
\]
Thus \(\mathcal C_T\), rather than \(\mathcal C_T^{\rm mod}\), is the
final PE correspondence certificate; \(M_T\) is the maximum witness
returned with it; and \(\mathcal B_T,K\) are the returned base and
hierarchy.  The last output of \noderef{HierarchySplitter} says that
\(K\) has the same core, PE closure, divider count, ordered stages, and
divider positions as \(H\).  We neither identify the two existential
correspondence certificates nor transport objects between them.

If the kind is SE or RE, set
\(\mathcal C_T=\mathcal C_T^{\rm mod}\) and \(M_T=M_T^{\rm mod}\).
\noderef{CorrespondingExtendedTashkinovTree} identifies the rung, every
positive ladder tree, every raw stored set, every exact connecting-edge
set, and every extension kind.  Define the new base core to be the
replacement terminal ladder tree.  The old nonparallel clause in
\noderef{HierarchyBase} identifies its old counterpart with \(T_n\).
Copy the support position and unused closure.  The terminal missing-color
field is reflexive for \(\sigma_n\), the parallel implication is vacuous,
and the nonparallel field is the defining core equality.  Copy the
divider count, stages, and divider positions of \(H\).  Every nestedness,
strictness, divider-prefix, and terminal-stage field of
\noderef{TashkinovHierarchy} is a coloring-free assertion about these
same ordered trees, so this defines \(\mathcal B_T\) and \(K\) field by
field.  In both terminal-kind branches, therefore, the objects in (9)
are now fixed and satisfy the existential correspondence, maximum,
base, and hierarchy requirements of the conclusion.

We compare the hierarchy colors before rebuilding goodness.
Correspondence identifies the raw stored arrays.  By
\noderef{TashkinovStoredColors}, the computed \(D_n\) removes from their
common union precisely the colors missing on \(T_n\).  Weak stability in
(S1) gives vertexwise equality of missing-color sets on \(T_n\), so the
computed terminal stored sets agree.  Every starred stage is an exact
prefix of \(A\preccurlyeq B\).  If \(v\) is a vertex of such a stage,
expand \noderef{TreeSequenceVertexSet} to choose its displayed index and
carry that index through the exact prefixes to \(B\); the third transport
hypothesis applies there.  Thus the two colorings have the same
tree-missing colors on every starred stage.  Expanding
\noderef{HierarchyStoredColorsAt} gives
\[
                         D^K_{n,j}=D^H_{n,j}                  \tag{10}
\]
at every stored index \(j\).  Expanding
\noderef{HierarchyStageMissingColors}, using \(T_n\) at index zero and
the retained positive stages thereafter, gives equality of all
stage-missing sets as well.                                      \(\tag{11}\)

Define the new Gamma pairs by copying
\(C.\mathit{good}.\Gamma^j_\eta\), using the canonical casts from the
equal divider counts.  Equations (10)--(11) transport the two-color and
pairwise-disjointness fields of \noderef{GoodHierarchy}.

Suppose \(q>0\).  The actual prefix witness \(g_A\) in (4a) supplies every
unique-refresh and minus-closed clause strictly before the last divider:
its stages, stored sets, and Gamma pairs are exactly the retained ones,
and all these objects are also the corresponding initial objects of
\(K\).  The separate last minus-closed conclusion of (S3) is exactly the
remaining clause for \(K\).  For the last unique-refresh clause, the
quantified stage difference is contained in \(A\preccurlyeq B\); the
third transport hypothesis, (10), and the copied Gamma pairs transport
the old unique witness and its uniqueness.  When \(q=0\), there are no
unique-refresh or minus-closed fields, and the parallel-base implication
is vacuous in the SE/RE case.

We next verify the PE parallel-base field, separating the two divider
cases.  First let \(q>0\).  The kind equality
\(\kappa_A=\kappa\) retained in (4a), together with the present assumption
\(\kappa=PE\), makes the antecedent of \(g_A.\mathit{parallelBase}\)
true.  Hence \(g_A\) gives both
\[
 \neg\noderef{TreeColorSetMinusClosed}
   (G,\sigma_n,\mathcal B_A.\mathit{core},\Gamma_A^0)           \tag{12}
\]
and
\[
 4\leq\bigl|(\overline{\sigma}_n(T_n^A)\cap
       \overline{\sigma}_n(\mathcal B_A.\mathit{closure}))
       -\Gamma_A^0\bigr|,                                      \tag{13}
\]
where \(T_n^A\) is the terminal ladder tree of \(\mathcal E_A\) and
\(\Gamma_A^0\) is the union of the zeroth \(g_A\)-pairs over the zeroth
stored set of \(H_A\).

Every set in (12)--(13) is exactly the corresponding set for the full
hierarchy \(K\).  Indeed, the core and closure equalities in (S3) identify
\(\mathcal B_A\) with the original base on those two fields, while the
PE output of \noderef{HierarchySplitter} identifies
\(\mathcal B_T\) with the same original core and closure.  Next,
the base-stage equation of \(H_A\), its retained stage-zero equality
with \(H\), the stage-zero equality returned for \(K\), and the
base-stage equation of \(K\) identify \(T_n^A\) with the terminal ladder
tree of \(\mathcal C_T.\mathit{replacementETT}\).  Finally the retained
zeroth stored-set and Gamma-pair equalities in (S3), equation (10), and
the copied definition of the new Gamma pairs identify
\(\Gamma_A^0\) with the zeroth union for \(K\).  Rewriting (12) and (13)
along these exact equalities proves, respectively, the non-minus-closed
and cardinality-at-least-four clauses for \(K\).  No edge-color transport
is involved: both assertions already hold under \(\sigma_n\).

It remains to verify the PE parallel-base clause when \(q=0\).  Its
cardinality assertion involves only the missing-color sets on \(T_n\) and
the PE closure, the stored set in (10), and the copied Gamma pairs.  The
two vertex clauses of the \noderef{CombinedStableColoring} relation in
(S1) identify these finite sets, so the old cardinality-at-least-four
assertion transports.

For non-minus-closedness, (2) makes (7) nonreflexive.  Take its first
\(\sigma_n\)-\noderef{TashkinovAugmentation} \(A\to Y\).  It supplies the
exact new edge \(g=Y.\mathit{edge}(|A|)\) and a color \(c\) such that
\[
                 \sigma_n(g)=c,\qquad
                 c\in\noderef{TreeMissingColors}(\sigma_n,A). \tag{14}
\]
The attachment and vertex-injectivity fields of
\noderef{TreeSequence}, together with the exact one-step prefix, show
that \(g\) has one endpoint in \(V(A)\) and the new endpoint outside it.
Thus \(g\) is a \noderef{BoundaryEdges} identity of \(V(A)\), and (14)
makes it a \noderef{BoundaryColorEdges} identity of color \(c\).

We claim \(c\notin\Gamma^0\).  If \(c=\alpha\) or \(c=\beta\), this is
the second starred-stage hypothesis.  Otherwise the definition of
\noderef{TwoColorComponentSwap} fixes every occurrence of the unaffected
color \(c\), in both directions, so (14) gives
\(\varphi_n(g)=c\).  The third transport hypothesis, restricted along
\(A\preccurlyeq B\), also turns the missing witness in (14) into a
\(\varphi_n\)-missing witness on \(A\).  If \(c\) belonged to
\(C.\mathit{good}.\Gamma^0_\eta\) for some
\(\eta\in D_{n,0}\), then
\noderef{HierarchyStoredColorsAt} says that \(\eta\) is not missing on
\(A\).  Hence its old
\noderef{HierarchyFirstMissingPosition} is strictly after \(|A|\), or is
the terminal fallback \(|T|\), which is also strictly after \(|A|\) by
(2).  The exact position \(|A|\), the edge \(g\), and
\(\varphi_n(g)=c\) would put \(c\) in
\noderef{HierarchySegmentColors}(H,0,\eta)\), contradicting the old
avoidance field of \noderef{GoodHierarchy}.  Thus \(c\notin\Gamma^0\) in
all cases.

Unfold \noderef{TreeColorSetMinusClosed} and
\noderef{TreeColorSetClosed}.  The boundary identity \(g\), (14), and
\(c\notin\Gamma^0\) witness that \(A\) is not
\((\Gamma^0)^-\)-closed under \(\sigma_n\).  This finishes the
parallel-base field without transporting an old boundary witness and
without fixing any affected edge after \(A\).

Finally verify avoidance.  Suppose first that \(q>0\), and fix a retained
index \(j<q\), a color \(\eta\) in the stored set at that index, and a
copied color \(c\in g_A.\Gamma^j_\eta\).  Let
\[
 U=H_A.\mathit{stage}(j+1)=K.\mathit{stage}(j+1)
\]
be the common next stage.  The retained stage equations give exact
initial segments \(U\preccurlyeq A\preccurlyeq T\), and in particular
\(|U|\leq|A|\).  We claim the exact numerical cutoff equality
\[
 \min\{\noderef{HierarchyFirstMissingPosition}
          (\sigma_n,A,\eta),|U|\}
 =
 \min\{\noderef{HierarchyFirstMissingPosition}
          (\sigma_n,T,\eta),|U|\}.                              \tag{15}
\]
Here and below the displayed first-missing positions mean their numeric
values.

To prove (15), unfold \noderef{HierarchyFirstMissingPosition}.  If
\(\eta\) has a missing witness on \(A\), let \(m_A\) be its least such
displayed index.  Entrywise vertex equality for the exact prefix
\(A\preccurlyeq T\) makes \(m_A\) a witness on \(T\), so the least
witness \(m_T\) on \(T\) satisfies \(m_T\leq m_A\).  This puts \(m_T\)
within the displayed range of \(A\); the same entrywise equality makes
it a witness on \(A\), and minimality of \(m_A\) gives
\(m_A\leq m_T\).  Thus \(m_A=m_T\), and (15) follows.

If \(\eta\) has no missing witness on \(A\), its totalized value on
\(A\) is the fallback \(|A|\).  If it has a witness on \(T\), the least
one has index at least \(|A|\), since any smaller witness would transport
back along \(A\preccurlyeq T\); if it has no witness on \(T\), its value
is the fallback \(|T|\geq|A|\).  In either subcase the right-hand
first-missing value is at least \(|A|\).  Since \(|U|\leq|A|\), both
minima in (15) equal \(|U|\).  This proves (15), including both fallback
possibilities.

The starred stage at \(j\), the next stage \(U\), and every exact edge of
\(U\) are the same for \(H_A\) and \(K\).  Equation (15) makes the upper
range bound in \noderef{HierarchySegmentColors} the same as well, and
both segment predicates use the same coloring \(\sigma_n\).  Therefore
the pointwise avoidance clause of \(g_A\) transfers to \(K\): its copied
color \(c\) is missing on the common stage specified by
\noderef{HierarchyStageMissingColors}, and no position in the common
segment has value \(c\).  This proves avoidance for every \(j<q\); the
fourth starred-stage hypothesis has not been used.

It remains to prove avoidance at \(j=q\), including the case \(q=0\).
Fix \(\eta\in D_{n,q}\) and
\(c\in C.\mathit{good}.\Gamma^q_\eta\).  Membership in the Gamma union
and the second transport hypothesis imply
\(c\notin\{\alpha,\beta\}\).  Therefore the component swap fixes every
edge whose old or new value is \(c\).  Old avoidance says that \(c\) is
stage-missing and is absent from
\noderef{HierarchySegmentColors}(H,q,\eta)\).  The stage-missing
membership transports by (11).  The fourth transport hypothesis makes
the first-missing endpoint in the definition of the old and new segment
identical.  Hence any exact position witnessing \(c\) in the new segment
would have the same range and the same edge of old color \(c\), contrary
to old avoidance.  This proves the last avoidance clause pointwise; no
generic equality of whole segment-color sets is asserted.

All fields of a \noderef{GoodHierarchy} \(g_K\) on \(K\) have now been
constructed.  The replacement base has the old core and PE closure when
applicable; the hierarchy has the same divider count, stages, and
positions; and (10) together with the copied definition of the Gamma
pairs gives equality of every stored set and pair.  These are exactly the
clauses of
\(\noderef{SameGoodHierarchy}(C.\mathit{good},g_K)\).  Combining this
with (S1), (S2), and the objects in (9) proves the lemma.
\end{proof}
